% GENERATED FILE. Edit canonical lessons, then run npm run generate:books.
% canonical-source-hash: fnv1a64:468a6ea8acc73400
% canonical-lessons: TA-C13-udal-uruppugal, TA-S06-ka, TA-W04-i-sign-write-nandri, TA-S117-letter-ra, TA-C13-checkpoint

\chapter{Body Parts}
\label{ch:ta-body-parts}
\hlchaptermodality{13}

\begin{chapteropening}
\textbf{By the end of this chapter:} \emph{I can say head and hand in Tamil, account for why both stayed native here where Hindi's pair came from Sanskrit, and write \ta{நன்றி} from its pieces.}

Complete TA-C13-checkpoint, the chapter's closing checkpoint: say talai and kai as native words, explain the d in nandri, describe \ta{க} and where \ta{◌ி} attaches, find \ta{ர} in \ta{சரி}, then write \ta{க}, \ta{ர}, \ta{றி} and \ta{நன்றி} from sound.
\end{chapteropening}

\section[\texorpdfstring{\ta{தலை} \ta{கை} (talai kai)}{தலை கை (talai kai)}]{\ta{தலை}, \ta{கை} — head and hand, both Tamil's own}

\label{lesson:TA-C13-udal-uruppugal}



\begin{quote}
You've seen Hindi's \emph{sir} and \emph{hāth} trace all the way back to Sanskrit. Tamil's words for the same two body parts tell a simpler story: \textbf{both are native.}
\end{quote}

\subsection*{You'll want to know: The words}
\begin{itemize}
\raggedright
  \item \textbf{\ta{தலை}} (\emph{talai}) = \textquotedblleft{}\textbf{head}\textquotedblright{} — native Dravidian.
  \item \textbf{\ta{கை}} (\emph{kai}) = \textquotedblleft{}\textbf{hand}\textquotedblright{} — native Dravidian.
\end{itemize}

Neither is borrowed from Sanskrit, unlike Hindi's \emph{sir} (< Sanskrit \emph{śiras}) or \emph{hāth} (< Sanskrit \emph{hasta}). Tamil simply kept its own words for two of the most basic parts of the body.

\textbf{In the family, and next door}

\noindent
\begin{tabularx}{\linewidth}{@{}>{\raggedright\arraybackslash}X>{\raggedright\arraybackslash}X>{\raggedright\arraybackslash}X@{}}
\toprule
\textbf{} & \textbf{word} & \textbf{same root} \\
\midrule
Kannada & \ka{ತಲೆ} \ka{ಕೈ} (\emph{tale kai}) & yes \\
Telugu & \te{తల} \te{చెయ్యి} (\emph{tala ceyyi}) & yes \\
Malayalam & \ml{തല} \ml{കൈ} (\emph{tala kai}) & yes \\
English & head and hand &  \\
\bottomrule
\end{tabularx}

\subsection*{Your turn}
Say these aloud:

\begin{itemize}
\raggedright
  \item \textquotedblleft{}talai\textquotedblright{} — head
  \item \textquotedblleft{}kai\textquotedblright{} — hand
  \item the contrast — Hindi borrowed these from Sanskrit; Tamil didn't
\end{itemize}

\begin{tcolorbox}[breakable,colback=teal!4,colframe=teal!35!black,title={Before you move on}]
What are Tamil's words for head and hand? (\textbf{\ta{தலை}} \emph{talai}, \textbf{\ta{கை}} \emph{kai}.) Are either of these Sanskrit loans? (\textbf{No} — both native Dravidian, unlike Hindi's Sanskrit-descended \emph{sir} and \emph{hāth}.)
\end{tcolorbox}
\section[\texorpdfstring{\ta{க} (ka)}{க (ka)}]{\ta{க} — the letter that is three bowls}

\label{lesson:TA-S06-ka}



\begin{quote}
Before the new one: say the sound of \textbf{\ta{ற}}.

One letter this time. Just one.

The most-used consonant in Tamil, and the most built: a square frame on top, then two bowls hanging under it. One continuous stroke, and a long one. Take it slowly — the pen climbs, crosses, drops and goes round both bowls before it lifts.
\end{quote}

\subsection*{Script you'll notice: \ta{க}}
\textbf{\ta{க}} — \emph{ka}.

What it is made of:

\begin{itemize}
\raggedright
  \item an upper square frame
  \item a large lower-left bowl
  \item a lower-right bowl with an open foot
\end{itemize}

\ta{வணக்கம்} has it twice, doubled in the middle.

\subsection*{Writing: \ta{க}}
\hlblockfigure{\detokenize{figures/TA-S06-ka-filmstrip.pdf}}{How \ta{க} is written, stroke by stroke}

\begin{itemize}
\raggedright
  \item \textbf{1.} start at the middle left and climb the left upright
  \item \textbf{2.} without lifting, carry the top bar to the right and return along it to the inner corner
  \item \textbf{3.} without lifting, drop the inner upright to the inner crossing
  \item \textbf{4.} without lifting, curve down and around the lower-left bowl
  \item \textbf{5.} without lifting, return up its outer left side to the middle left
  \item \textbf{6.} without lifting, cross the middle bar and turn around the lower-right bowl — and only now lift
\end{itemize}

\textbf{Pen lifts: 0.} The pen never leaves the paper.

\begin{quote}
Stroke order is one attested teaching order, not a national standard — Tamil handwriting is taught with school-to-school variation. Source: Sankaran Radhakrishnan, Tamil Script Learners Manual, Appendix I: Hand-movements, Frame 3, \ta{க} (Univ. of Texas at Austin), p. 191.
\end{quote}

\subsection*{Your turn}
\begin{itemize}
\raggedright
  \item \emph{Trace it:} \ta{க} once, slowly, following the steps above
  \item \emph{Say it:} \textquotedblleft{}ka\textquotedblright{} as you finish it
  \item \emph{Trace it:} it twice more, without looking back at the steps
\end{itemize}

\begin{tcolorbox}[breakable,colback=teal!4,colframe=teal!35!black,title={Before you move on}]
How many pen lifts does \ta{க} take? (\textbf{None} — it is one continuous stroke.) What are the three parts? (\textbf{Upper frame, lower-left bowl, lower-right bowl.})
\end{tcolorbox}
\section[\texorpdfstring{\ta{ி}, \ta{நன்றி} (i, naṉṟi / nandri)}{ி, நன்றி (i, naṉṟi / nandri)}]{\ta{ி} and \ta{நன்றி} — replace the vowel, then hear \emph{ndr}}

\label{lesson:TA-W04-i-sign-write-nandri}



\begin{quote}
A consonant can keep \emph{a} or lose it under a puḷḷi. The third option is to replace it with a vowel sign.
\end{quote}

\subsection*{Script you'll notice: The first vowel sign}
\hlblockfigure{\detokenize{figures/TA-W04-i-sign-write-nandri-filmstrip.pdf}}{How \ta{ி}, \ta{நன்றி} is written, part by part, stroke by stroke}

\begin{quote}
\textbf{\ta{ற}} = \emph{ṟa} $\to$ \textbf{\ta{றி}} = \emph{ṟi}
\end{quote}

\textbf{\ta{ி}} is a hook above and to the right. The consonant keeps its shape; the sign hangs from it. Later, \textbf{\ta{ை}} (\emph{ai}) will appear before its consonant on the page but after it in speech, the same eye/ear trap as Devanagari \dv{ि}.

\subsection*{Script you'll notice: Assemble \ta{நன்றி}}
\begin{quote}
\textbf{\ta{ந} · \ta{ன்} · \ta{றி} $\to$ \ta{நன்றி}}
\end{quote}

\noindent
\begin{tabularx}{\linewidth}{@{}>{\raggedright\arraybackslash}X>{\raggedright\arraybackslash}X>{\raggedright\arraybackslash}X@{}}
\toprule
\textbf{piece} & \textbf{} & \textbf{} \\
\midrule
\textbf{\ta{ந}} & \emph{na} — dental & last lesson \\
\textbf{\ta{ன்}} & \emph{ṉ} — alveolar, with puḷḷi & Lessons 3–4 \\
\textbf{\ta{றி}} & \emph{ṟi} — \ta{ற} plus the i-sign & this lesson \\
\bottomrule
\end{tabularx}

\begin{sounds}
\textbf{\ta{ன்} + \ta{ற}} is pronounced close to \textbf{ndr}. This follows a broader Tamil rule: a \textbf{nasal voices the stop after it}.

\noindent
\begin{tabularx}{\linewidth}{@{}>{\raggedright\arraybackslash}X>{\raggedright\arraybackslash}X>{\raggedright\arraybackslash}X@{}}
\toprule
\textbf{} & \textbf{} & \textbf{} \\
\midrule
\ta{ந்} + \ta{த} & \emph{nd} & \textbf{\ta{வந்து}} \emph{vandu} \\
\ta{ண்} + \ta{ட} & \emph{ṇḍ} & retroflex pair \\
\ta{ன்} + \ta{ற} & \emph{ndr} & \textbf{\ta{நன்றி}} \\
\bottomrule
\end{tabularx}

Each n-letter creates its own cluster. That payoff explains both the three spellings and the common English romanization \textbf{nandri}, whose \emph{d} is heard but not separately written.
\end{sounds}

\subsection*{Your turn}
\begin{itemize}
\raggedright
  \item \emph{Write it:} \ta{றி} — \ta{ற} plus the hook above right
  \item \emph{Write it:} \textbf{\ta{நன்றி}} — three pieces
  \item \emph{Say it:} \textquotedblleft{}\ta{ன்} + \ta{ற} = \textbf{ndr}\textquotedblright{}
\end{itemize}

\begin{tcolorbox}[breakable,colback=teal!4,colframe=teal!35!black,title={Before you move on}]
What are a consonant's three vowel choices? (Keep inherent \textbf{a}, \textbf{remove} it with puḷḷi, or \textbf{replace} it with a vowel sign.) Where does \textbf{\ta{ி}} attach? (\textbf{Above and right}.) Why is \emph{naṉṟi} often written \emph{nandri}? (\textbf{\ta{ன்} + \ta{ற}} is pronounced \emph{ndr} because a nasal voices the following stop.)
\end{tcolorbox}
\section[\texorpdfstring{\ta{ர} (ra)}{ர (ra)}]{\ta{ர} — one character, met inside words you already say}

\label{lesson:TA-S117-letter-ra}



\begin{quote}
Before the new one: \ta{ச} — what does it do?

One character this time. Just one — and you have been saying it for pages without knowing which mark on the page it was.
\end{quote}

\subsection*{Script you'll notice: \ta{ர}}
\textbf{\ta{ர}} — \emph{ra}.

It is a \textbf{consonant}, and in this script a consonant is never bare: it comes with an \emph{a} already in it. So it is not \emph{r}, it is \textbf{ra}.

You already say these, and every one of them has \ta{ர} somewhere inside it:

\begin{itemize}
\raggedright
  \item \textbf{\ta{வணக்கம்} / \ta{நமஸ்காரம்}} \emph{vaṇakkam / namaskāram} — Tamil kept a native bow-word where its neighbors borrowed Sanskrit
  \item \textbf{\ta{சரி}} \emph{sari} — okay / alright / correct (sari)
  \item \textbf{\ta{ஆம்} / \ta{இல்லை} / \ta{சரி}} \emph{ām / illai / sari} — answering in Tamil
  \item \textbf{\ta{பெயர்}} \emph{peyar} — name
\end{itemize}

\subsection*{Writing: \ta{ர} — follow the numbered strip}
\hlblockfigure{\detokenize{figures/TA-S117-letter-ra-filmstrip.pdf}}{How \ta{ர} is written, stroke by stroke}

Follow the numbered strip: start where movement 1 begins, and let each panel show you the next movement before your pen makes it. Then write \ta{ர} yourself — slowly, and larger than it is printed.

\begin{quote}
The strip shows \textbf{where to start the character and which way to travel}, in one attested order. That is taught with real variation from school to school, so the strip names its source beneath it: treat it as a sound way in, not the only one.
\end{quote}

\subsection*{Your turn}
\begin{itemize}
\raggedright
  \item \emph{Look:} at these words, and find \ta{ர} in the ones that have it
\end{itemize}

\begin{quote}
\ta{சரி}  ·  \ta{வணக்கம்}
\end{quote}

\begin{itemize}
\raggedright
  \item \emph{Trace it:} \ta{ர} three times, saying \emph{ra} as you finish each one
  \item \emph{Look:} back at any page of this chapter and find \ta{ர} once more
\end{itemize}

\begin{tcolorbox}[breakable,colback=teal!4,colframe=teal!35!black,title={Before you move on}]
Which character is this — \ta{ர}? What sound does it carry? (\emph{\textbf{ra}}.) Name one word you already say that contains it.
\end{tcolorbox}
\section[Practice]{Chapter 13 checkpoint — head, hand, and \ta{நன்றி} from sound}

\label{lesson:TA-C13-checkpoint}



\begin{quote}
Nothing new here. The two body words out loud first, then two letters, a vowel sign and the word \emph{thank you} on the page.
\end{quote}

\subsection*{Your turn — head and hand}
\begin{itemize}
\raggedright
  \item \emph{Say it:} \textquotedblleft{}head\textquotedblright{}, then \textquotedblleft{}hand\textquotedblright{}, in Tamil
\end{itemize}

(\emph{Talai}; \emph{kai}.)

\begin{itemize}
\raggedright
  \item Are either of them Sanskrit loans, the way Hindi's \emph{sir} and \emph{hāth} are? (No: both are native Dravidian.)
  \item Why do people spell \emph{thank you} as \emph{nandri} when no \emph{d} is written? (The nasal voices the stop after it: \emph{ṉ} plus \emph{ṟ} is said close to \emph{ndr}.)
\end{itemize}

\subsection*{Script — \ta{க}, \ta{ர} and the i-sign}
\begin{itemize}
\raggedright
  \item What are the three parts of \textbf{\ta{க}}, and how many times does the pen lift? (An upper frame, a lower-left bowl, a lower-right bowl; no lift at all.)
  \item Where does \textbf{\ta{ி}} attach, and what does it do to the consonant's \emph{a}? (Above and to the right; it replaces the \emph{a}, so \textbf{\ta{ற}} becomes \textbf{\ta{றி}}.)
  \item Which three pieces build \textbf{\ta{நன்றி}}? (\textbf{\ta{ந}} · \textbf{\ta{ன்}} · \textbf{\ta{றி}}.)
  \item \textbf{\ta{சரி}} is \emph{okay}. Which letter in it carries \emph{ra}? (\textbf{\ta{ர}}.)
\end{itemize}

\emph{Read it:} \textbf{\ta{நன்றி}} and \textbf{\ta{சரி}} aloud, and point to the \ta{ர}

\subsection*{Writing — from sound}
\emph{Cover:} the Script block above, so no model stays in view

\emph{Write it:} the letter for \emph{ka}, then the letter for \emph{ra}

\emph{Write it:} \emph{ṟi}, then the whole of \emph{naṉṟi}, \textquotedblleft{}thank you\textquotedblright{}

\emph{Check:} against the key — \textbf{\ta{க}}, \textbf{\ta{ர}}, \textbf{\ta{றி}}, \textbf{\ta{நன்றி}}

\emph{Write it:} one repaired shape, then stop

\begin{tcolorbox}[breakable,colback=teal!4,colframe=teal!35!black,title={Before you move on}]
Touch your head and say its word, then your hand and say its word. That is the chapter: two native words, and \emph{thank you} written from its pieces.

Next chapter: the seasons.
\end{tcolorbox}
